% GENERATED FILE. Edit canonical lessons, then run npm run generate:books.
% canonical-source-hash: fnv1a64:2ee6bf36a48c5913
% canonical-lessons: JA-C529-nazeka, JA-C529-tokorode, JA-C529-douse, JA-C529-isso, JA-C529-sasuga

\chapter{Anyway and After All}
\label{ch:ja-anyway-and-after-all}
\hlchaptermodality{529}

\begin{chapteropening}
\textbf{By the end of this chapter:} \emph{I can comment on how things will turn out.}

Complete the last lesson of chapter 529: I can comment on how things will turn out.
\end{chapteropening}

\section[\texorpdfstring{\ja{なぜか} (nazeka)}{なぜか (nazeka)}]{\ja{なぜか} (nazeka) — for some reason}

\label{lesson:JA-C529-nazeka}



\begin{quote}
Before the new one: say the Japanese for or else, then the Japanese for either one.
\end{quote}

\subsection*{You'll want to know: \ja{なぜか}}
\textbf{\ja{なぜか}} — \emph{nazeka} — \textquotedblleft{}for some reason\textquotedblright{}.

\begin{quote}
\textbf{\ja{なぜか} \ja{あじは} \ja{みずに} \ja{よります。}} — \emph{nazeka aji wa mizu ni yorimasu} — \textquotedblleft{}For some reason, the taste depends on the water.\textquotedblright{}
\end{quote}

\begin{cousinweb}
\textbf{\ja{なぜか}} (\emph{nazeka}) is a native Japanese word (wago), inherited rather than borrowed from Chinese, and usually written in kana: built from \emph{naze}, \textquotedblleft{}why\textquotedblright{}, and the particle \emph{ka} (JMdict).
\end{cousinweb}

\begin{grammarlens}[title={What you've built}]
One more word for saying what happens if.
\end{grammarlens}

\subsection*{Your turn}
\begin{itemize}
\raggedright
  \item \emph{Say it:} \emph{nazeka}
  \item \emph{Say it:} \emph{nazeka}, once more
  \item \emph{Say it:} \emph{nazeka aji wa mizu ni yorimasu} — For some reason, the taste depends on the water.
  \item \emph{Recall:} say the Japanese for or else, then the Japanese for either one, then say \emph{nazeka} again
\end{itemize}

\begin{tcolorbox}[breakable,colback=teal!4,colframe=teal!35!black,title={Before you move on}]
What does \ja{なぜか} mean? (\textquotedblleft{}For some reason\textquotedblright{}.) Say it once more.
\end{tcolorbox}
\section[\texorpdfstring{\ja{ところで} (tokorode)}{ところで (tokorode)}]{\ja{ところで} (tokorode) — by the way}

\label{lesson:JA-C529-tokorode}



\begin{quote}
Before the new one: say the Japanese for either one, then the Japanese for for some reason.
\end{quote}

\subsection*{You'll want to know: \ja{ところで}}
\textbf{\ja{ところで}} — \emph{tokorode} — \textquotedblleft{}by the way\textquotedblright{}.

\begin{quote}
\textbf{\ja{ところで、もし} \ja{じかんが} \ja{あったら、えいがに} \ja{いきませんか。}} — \emph{tokorode, moshi jikan ga attara, eiga ni ikimasen ka} — \textquotedblleft{}By the way, if you have time, would you like to see a film?\textquotedblright{}
\end{quote}

\begin{cousinweb}
\textbf{\ja{ところで}} (\emph{tokorode}) is a native Japanese word (wago), inherited rather than borrowed from Chinese, and usually written in kana: built from \emph{tokoro}, \textquotedblleft{}place\textquotedblright{}, and the particle \emph{de} (JMdict).
\end{cousinweb}

\begin{grammarlens}[title={What you've built}]
One more word for saying what happens if.
\end{grammarlens}

\subsection*{Your turn}
\begin{itemize}
\raggedright
  \item \emph{Say it:} \emph{tokorode}
  \item \emph{Say it:} \emph{tokorode}, once more
  \item \emph{Say it:} \emph{tokorode, moshi jikan ga attara, eiga ni ikimasen ka} — By the way, if you have time, would you like to see a film?
  \item \emph{Recall:} say the Japanese for either one, then the Japanese for for some reason, then say \emph{tokorode} again
\end{itemize}

\begin{tcolorbox}[breakable,colback=teal!4,colframe=teal!35!black,title={Before you move on}]
What does \ja{ところで} mean? (\textquotedblleft{}By the way\textquotedblright{}.) Say it once more.
\end{tcolorbox}
\section[\texorpdfstring{\ja{どうせ} (dōse)}{どうせ (dōse)}]{\ja{どうせ} (dōse) — anyway; in any case}

\label{lesson:JA-C529-douse}



\begin{quote}
Before the new one: say the Japanese for for some reason, then the Japanese for by the way.
\end{quote}

\subsection*{You'll want to know: \ja{どうせ}}
\textbf{\ja{どうせ}} — \emph{dōse} — \textquotedblleft{}anyway; in any case\textquotedblright{}.

\begin{quote}
\textbf{\ja{もし} \ja{あめが} \ja{ふったら、どうせ} \ja{しあいは} \ja{ちゅうしでしょう。}} — \emph{moshi ame ga futtara, dōse shiai wa chūshi deshō} — \textquotedblleft{}If it rains, the match will be called off anyway.\textquotedblright{}
\end{quote}

\begin{cousinweb}
\textbf{\ja{どうせ}} (\emph{dōse}) is usually counted a native Japanese word (wago), inherited rather than borrowed from Chinese: JMdict gives it no kanji spelling, and no dictionary we use records an older form, so its history is not settled (JMdict).
\end{cousinweb}

\begin{grammarlens}[title={What you've built}]
One more word for saying what happens if.
\end{grammarlens}

\subsection*{Your turn}
\begin{itemize}
\raggedright
  \item \emph{Say it:} \emph{dōse}
  \item \emph{Say it:} \emph{dōse}, once more
  \item \emph{Say it:} \emph{moshi ame ga futtara, dōse shiai wa chūshi deshō} — If it rains, the match will be called off anyway.
  \item \emph{Recall:} say the Japanese for for some reason, then the Japanese for by the way, then say \emph{dōse} again
\end{itemize}

\begin{tcolorbox}[breakable,colback=teal!4,colframe=teal!35!black,title={Before you move on}]
What does \ja{どうせ} mean? (\textquotedblleft{}Anyway; in any case\textquotedblright{}.) Say it once more.
\end{tcolorbox}
\section[\texorpdfstring{\ja{いっそ} (isso)}{いっそ (isso)}]{\ja{いっそ} (isso) — would sooner; rather}

\label{lesson:JA-C529-isso}



\begin{quote}
Before the new one: say the Japanese for by the way, then the Japanese for anyway.
\end{quote}

\subsection*{You'll want to know: \ja{いっそ}}
\textbf{\ja{いっそ}} — \emph{isso} — \textquotedblleft{}would sooner; rather\textquotedblright{}.

\begin{quote}
\textbf{\ja{もし} \ja{あしたも} \ja{まけたら、いっそ} \ja{やめた} \ja{ほうが} \ja{いいです。}} — \emph{moshi ashita mo maketara, isso yameta hō ga ii desu} — \textquotedblleft{}If we lose tomorrow as well, it would be better to quit altogether.\textquotedblright{}
\end{quote}

\begin{cousinweb}
\textbf{\ja{いっそ}} (\emph{isso}) is usually counted a native Japanese word (wago), inherited rather than borrowed from Chinese, though that is not certain: JMdict gives it no kanji spelling at all, and no dictionary we use records its history, so where it comes from is not settled (JMdict).
\end{cousinweb}

\begin{grammarlens}[title={What you've built}]
One more word for saying what happens if.
\end{grammarlens}

\subsection*{Your turn}
\begin{itemize}
\raggedright
  \item \emph{Say it:} \emph{isso}
  \item \emph{Say it:} \emph{isso}, once more
  \item \emph{Say it:} \emph{moshi ashita mo maketara, isso yameta hō ga ii desu} — If we lose tomorrow as well, it would be better to quit altogether.
  \item \emph{Recall:} say the Japanese for by the way, then the Japanese for anyway, then say \emph{isso} again
\end{itemize}

\begin{tcolorbox}[breakable,colback=teal!4,colframe=teal!35!black,title={Before you move on}]
What does \ja{いっそ} mean? (\textquotedblleft{}Would sooner; rather\textquotedblright{}.) Say it once more.
\end{tcolorbox}
\section[\texorpdfstring{\ja{さすが} (sasuga)}{さすが (sasuga)}]{\ja{さすが} (sasuga) — as one would expect}

\label{lesson:JA-C529-sasuga}



\begin{quote}
Before the new one: say the Japanese for anyway, then the Japanese for would sooner.
\end{quote}

\subsection*{You'll want to know: \ja{さすが}}
\textbf{\ja{さすが}} — \emph{sasuga} — \textquotedblleft{}as one would expect\textquotedblright{}.

\begin{quote}
\textbf{\ja{もし} \ja{あの} \ja{せんしゅが} \ja{かったら、さすがでしょう。}} — \emph{moshi ano senshu ga kattara, sasuga deshō} — \textquotedblleft{}If that athlete wins, it will be what you would expect of her.\textquotedblright{}
\end{quote}

\begin{cousinweb}
\textbf{\ja{さすが}} (\emph{sasuga}) is a native Japanese word (wago), inherited rather than borrowed from Chinese, and usually written in kana; its older origin is not settled (JMdict).
\end{cousinweb}

\begin{grammarlens}[title={What you've built}]
One more word for saying what happens if.
\end{grammarlens}

\subsection*{Your turn}
\begin{itemize}
\raggedright
  \item \emph{Say it:} \emph{sasuga}
  \item \emph{Say it:} \emph{sasuga}, once more
  \item \emph{Say it:} \emph{moshi ano senshu ga kattara, sasuga deshō} — If that athlete wins, it will be what you would expect of her.
  \item \emph{Recall:} say the Japanese for anyway, then the Japanese for would sooner, then say \emph{sasuga} again
\end{itemize}

\begin{tcolorbox}[breakable,colback=teal!4,colframe=teal!35!black,title={Before you move on}]
What does \ja{さすが} mean? (\textquotedblleft{}As one would expect\textquotedblright{}.) Say it once more.
\end{tcolorbox}
